\documentclass[11pt,a4paper]{article}

% ===== XeLaTeX + Korean =====
\usepackage{kotex}
\usepackage{fontspec}
\setmainfont{Times New Roman}
\setmainhangulfont{AppleMyungjo}
\setsanshangulfont{Apple SD Gothic Neo}

% ===== Layout =====
\usepackage[a4paper,top=18mm,bottom=18mm,left=20mm,right=20mm]{geometry}
\usepackage{fancyhdr}
\setlength{\headheight}{24pt}
\usepackage{enumitem}
\usepackage{mdframed}
\usepackage{array}

\setlength{\parindent}{1em}
\linespread{1.18}

% ===== Header / Footer =====
\pagestyle{fancy}
\fancyhf{}
\renewcommand{\headrulewidth}{0.6pt}
\fancyhead[L]{\sffamily\bfseries 2026 고1 영어 변형 — 지문 24}
\fancyhead[R]{\sffamily 도파민과 SNS}
\renewcommand{\footrulewidth}{0pt}
\fancyfoot[C]{\framebox[16mm][c]{\thepage}}

% ===== helpers =====
\newcommand{\qno}[1]{\par\vspace{8pt}\noindent{\sffamily\bfseries #1.}\ }
\newcommand{\instr}[1]{{\sffamily #1}\par\vspace{3pt}}
\newlist{choices}{enumerate}{1}
\setlist[choices]{label=,leftmargin=1.5em,itemsep=2pt,topsep=3pt,parsep=0pt}

\newmdenv[linewidth=0.6pt,roundcorner=3pt,innertopmargin=7pt,innerbottommargin=7pt,
          innerleftmargin=9pt,innerrightmargin=9pt,skipabove=5pt,skipbelow=5pt]{pbox}
\newcommand{\nemo}[1]{\fbox{\,#1\,}}

\begin{document}

\begin{center}
{\fontsize{15}{17}\selectfont\sffamily\bfseries [지문 24] 다음 글을 읽고 물음에 답하시오.}
\end{center}
\vspace{4pt}

% ===== 공통 지문 (homeostasis 첫 번째 → 유지, 두 번째 → 빈칸 / seeking → avoiding (어휘 오답)) =====
\begin{pbox}
Let's look at what is \underline{①occurring} in your brain when scrolling on social media, a quick dopamine activity. You open your social media apps and immediately you notice you feel really good, your dopamine levels increase \underline{②incredibly} fast and pleasure is experienced in your brain. The challenge this creates is as with everything in our universe: as the laws of physics explain, `what goes up, must come down'. Your brain and body are always \underline{③avoiding} something called `homeostasis', which simply means balance. With this in mind, when dopamine levels increase incredibly fast when scrolling social media, the brain then thinks, `Wow, how are my dopamine levels so high?' In order to achieve \rule{3.5cm}{0.4pt}, or balance, the dopamine then has to quickly drop an equal amount below your \underline{④baseline} level in order to \underline{⑤rebalance}, making you feel even worse than before you started scrolling.
\end{pbox}

% ===================== Q1 =====================
\qno{1}\instr{다음 빈칸에 들어갈 말로 가장 적절한 것은?}
\begin{choices}
\item ① stimulation
\item ② equilibrium
\item ③ acceleration
\item ④ gratification
\item ⑤ withdrawal
\end{choices}

% ===================== Q2 =====================
\qno{2}\instr{밑줄 친 ①$\sim$⑤ 중에서 문맥상 낱말의 쓰임이 적절하지 않은 것은?}
\begin{choices}
\item ① ①\quad ② ②\quad ③ ③\quad ④ ④\quad ⑤ ⑤
\end{choices}

% ===================== Q3 (서술형) =====================
\qno{3}\instr{[서술형] 다음 문장을 우리말로 해석하시오.}
\vspace{2pt}\par\noindent
\hspace{1em}\textit{the dopamine then has to quickly drop an equal amount below your baseline level in order to rebalance, making you feel even worse than before you started scrolling.}
\par\vspace{6pt}\noindent
$\Rightarrow$ \rule{\linewidth}{0.4pt}
\par\vspace{4pt}\noindent
\rule{\linewidth}{0.4pt}

% ===================== Q4 =====================
\qno{4}\instr{윗글의 제목으로 가장 적절한 것은?}
\begin{choices}
\item ① Social Media as the Safest Source of Lasting Pleasure
\item ② Why Dopamine Keeps Rising the Longer You Scroll
\item ③ The Physics of Social Networks and Online Communication
\item ④ Homeostasis: The Brain's Refusal to Feel Pleasure
\item ⑤ The Quick Dopamine High and the Cost of Rebalancing
\end{choices}

\end{document}
